% ==================================================================
% SU(3) Color Charge from Tetrahedral Cage Representations
% in Conscious Point Physics
% Companion Paper 15 to “Mechanistic Derivation of Relativistic Effects
% via Space Stress Vector (SSV) in the Dipole Sea” (Version 16)
% 19 March 2026 — Version 1 (independent Grok draft)
% ==================================================================

\documentclass[12pt]{article}

\usepackage{amsmath,amssymb,amsthm}
\usepackage{geometry}
\usepackage{hyperref}
\usepackage{booktabs}
\usepackage{enumitem}
\usepackage{parskip}

\geometry{letterpaper, margin=1in}

\newtheorem{theorem}{Theorem}[section]
\newtheorem{proposition}[theorem]{Proposition}
\newtheorem{corollary}[theorem]{Corollary}
\newtheorem{remark}[theorem]{Remark}

% ── CPP macros ────────────────────────────────────────────────────────────────
\newcommand{\lP}{l_P}
\newcommand{\SSV}{\rm SSV}
\newcommand{\qCP}{\rm qCP}
\newcommand{\qDP}{\rm qDP}
\newcommand{\as}{\alpha_s}
\newcommand{\sig}{\sigma}
\newcommand{\rzero}{r_0}

\title{\textbf{SU(3) Color Charge from Tetrahedral Cage Representations\\
in Conscious Point Physics}\\[6pt]
{\large Companion Paper 15 ---
Companion Paper 15 to ``Mechanistic Derivation of Relativistic Effects\\
via Space Stress Vector (SSV) in the Dipole Sea'' (Version~16)}}

\author{Thomas Lee Abshier, ND\\
Hyperphysics Institute\\
\texttt{https://hyperphysics.com}\\
\texttt{drthomas007@protonmail.com}}

\date{19 March 2026 --- Version~1 (independent Grok draft)}

\begin{document}
\maketitle

\noindent\textbf{Keywords:} Conscious Point Physics, SU(3) color, tetrahedral cage, 600-cell, Z₃ center, color triplet, gluon octet, color neutrality, qDP chain.

%======================================================================
\section*{Plain Language Summary}
%======================================================================

Quarks carry a hidden “color” charge that comes in three varieties (red, green, blue). Only color-neutral combinations (meson \(q\bar{q}\), baryon \(qqq\), glueball) are stable; free quarks are forbidden. In QCD this is an input from the SU(3) gauge group. In CPP it emerges geometrically from the tetrahedral cage that every quark Conscious Point occupies inside the 600-cell lattice.

A tetrahedron has one apex vertex and three base vertices. Rotation by 120° around the apex-to-centroid axis cycles the three base vertices — exactly the three color states. The 1+3+4 layered structure of the qDP chain between quarks maps onto the 8 gluon generators. Color neutrality is the algebraic requirement that the sum of color charges is zero modulo 3. All of this follows from the same cage geometry that already gave fractional electric charges (companion 2) and confinement (companion 14). No new postulates.

%======================================================================
\begin{abstract}
We derive the SU(3) color structure of quarks from the tetrahedral cage geometry of the 600-cell lattice. The qCP cage is a regular tetrahedron (Td symmetry). Its apex vertex plus three base vertices generate a color triplet under the C₃ rotation around the apex-centroid axis (Theorem 3.1). Color charge is the Z₃ quantum number carried by the chain-attachment point on the base triangle. The 1+3+4 layered qDP chain between quarks maps onto the 8 gluon generators: the central + constrained outer chains give the 2 diagonal Cartan generators, while the three middle-layer chains at 120° give the 6 off-diagonal raising/lowering operators (Proposition 4.1). Color neutrality requires the total Z₃ charge ≡ 0 mod 3, which algebraically selects only meson (\(q\bar{q}\)), baryon (\(qqq\)), and glueball configurations (Theorem 5.1). Fractional electric charges and color share the same tetrahedral projection origin (Remark 3.2). The result closes the strong-force sector: confinement mechanism (companion 14) + color algebra (this paper) together give QCD phenomenology from first principles.
\end{abstract}

%======================================================================
\section{Introduction}
%======================================================================

Companion 14 derived confinement and the Cornell potential from qDP chaining, but treated color charge as an input. This paper closes that gap: SU(3) color, the triplet, the octet, and the singlet condition emerge geometrically from the tetrahedral cage that every quark Conscious Point occupies inside the 600-cell lattice.

The cage is the same object that already produced fractional electric charges via 3D projection (companion 2). Here we show that the same tetrahedron also encodes the three color states via its discrete rotational symmetry and the 8 gluon modes via the layered qDP chain structure. No new postulates are introduced.

%======================================================================
\section{The Tetrahedral qCP Cage}
%======================================================================

Every quark Conscious Point occupies a regular tetrahedral cage inscribed in the Voronoi cell of the 600-cell lattice (Version 16). The tetrahedron has four vertices:
- one apex vertex (primary chain origin),
- three base vertices forming an equilateral triangle.

The symmetry group of the cage is Td (order 24). Crucially, there exists a C₃ rotation axis passing through the apex vertex and the centroid of the base triangle.

\begin{theorem}[Color triplet from C₃ rotation]
A 120° rotation around the apex-centroid axis cyclically permutes the three base vertices. These three vertices, related by the C₃ action, are the geometric origin of the three color states (red, green, blue).
\end{theorem}

The color state of a quark is defined by which base vertex carries the active qDP-chain attachment point. Rotation of the chain attachment point by 120° changes red → green → blue → red.

\begin{remark}[Shared origin with fractional charge]
The same tetrahedral cage already gave fractional electric charges via projection onto 3D space (companion 2). Color charge arises from the rotational degeneracy of the same three base vertices. The two properties — electric charge in thirds and three color states — are therefore two different views of one geometric object.
\end{remark}

%======================================================================
\section{Color Charge as Z₃ Quantum Number}
%======================================================================

The C₃ rotation generates the cyclic group Z₃. We assign color charge as a Z₃ quantum number (0, 1, 2 mod 3) to the active base vertex:
- vertex 1: color 0,
- vertex 2: color 1,
- vertex 3: color 2.

A color rotation is exactly a C₃ lattice operation that cycles the labels. The center of the color group is therefore Z₃ — recovered directly from the cage symmetry without postulating SU(3).

%======================================================================
\section{The Gluon Octet from qDP Chain Layers}
%======================================================================

From companion 14 the qDP chain between a quark and antiquark has three concentric layers:
- 1 central chain,
- 3 middle-layer chains at 120° intervals,
- 4 outer-layer chains (geometrically constrained by the tetrahedral symmetry).

These 1+3+4 = 8 modes map onto the 8 generators of SU(3) as follows:

\begin{proposition}[Gluon octet from chain layers]
- Central chain + one geometrically constrained outer chain → the 2 diagonal Cartan generators (T₃, Y).
- The three middle-layer chains, each able to act in two directions (raising/lowering the color label), → the 6 off-diagonal generators.
- Total: 2 + 6 = 8 generators.
\end{proposition}

A gluon is a transverse oscillation of the qDP chain that moves the attachment point on the base triangle — precisely a color rotation. The layer structure therefore reproduces the full adjoint representation counting of SU(3).

%======================================================================
\section{Color Neutrality as Algebraic Singlet Condition}
%======================================================================

\begin{theorem}[Color neutrality]
Only configurations whose total color charge sums to 0 mod 3 are stable. A free quark (open chain) carries nonzero Z₃ charge and requires an infinite-energy chain, which is forbidden (companion 14, Theorem 6.1).
\end{theorem}

Proof: meson (\(q\bar{q}\)): opposite colors sum to 0 mod 3; baryon (\(qqq\)): three identical colors sum to 0 mod 3; glueball: no qCP endpoints, net charge 0.

This algebraic condition is stronger than the topological closure argument of companion 14 and selects exactly the observed color-neutral states.

%======================================================================
\section{Consistency with Previous Companions}
%======================================================================

- Tetrahedral cage geometry and fractional-charge projection: companion 2 and Version 16.
- qDP chain layers (1+3+4): companion 14.
- Confinement and infinite-energy open chain: companion 14, Theorem 6.1.
- CP Exclusion and PSR: companion 1 (used implicitly in chain stability).

The color structure is derived from the same objects that already gave electromagnetism, gravity, spin, and confinement. No new postulates.

%======================================================================
\section{Open Problems}
%======================================================================

1. Full SU(3) Lie-algebra commutation relations from cage representations (beyond counting).
2. Quantitative derivation of the strong coupling running and β-function from PSR + chain self-energy (companion 14 Open Problem 2).
3. Embedding of the full Standard-Model gauge group into H₄ representations (electroweak sector).

%======================================================================
\section{Conclusion}
%======================================================================

1. Color triplet: C₃ rotation of the tetrahedral base vertices (Theorem 3.1).
2. Z₃ center and color charge: direct from cage symmetry.
3. Gluon octet: 1+3+4 qDP chain layers mapped to 8 generators (Proposition 4.1).
4. Color neutrality: algebraic sum ≡ 0 mod 3 (Theorem 5.1).
5. Shared origin with fractional electric charge: same tetrahedral cage.

SU(3) color is no longer an input. It is a geometric fact of the 600-cell lattice. Confinement (companion 14) + color algebra (this paper) together close the strong-force sector from first principles.

\textbf{GitHub:} 600-cell cage symmetry routines (extension of the existing Monte-Carlo repository) for future Lie-algebra verification.

\end{document}
